Neo4j to otwartoźródłowa grafowa baza danych rozwijana od 2007 roku przez
szwedzką firmę Neo4j Inc. Obecna linia 5.x (stan na 2025~r.) należy do
kategorii baz NoSQL i zamiast tabel czy dokumentów przechowuje dane w postaci
grafu złożonego z węzłów (nodes), krawędzi (relationships) i właściwości
(properties). Do odpytywania danych służy deklaratywny język Cypher,
zaprojektowany specjalnie do nawigacji po grafach. Neo4j jest dostępny
na licencji GPL (Community Edition) oraz komercyjnej (Enterprise Edition).

\subsection{Skalowalność}

Neo4j w wersji Community Edition działa jako pojedynczy węzeł bez możliwości
replikacji. Enterprise Edition oferuje Causal Clustering -- mechanizm
klastrowania oparty o protokół Raft, który zapewnia replikację i wysoką
dostępność. Klaster składa się z węzłów primary (obsługują zapis) i secondary
(obsługują odczyt), a spójność między nimi jest gwarantowana przez protokół
consensusu.

Skalowanie poziome przez sharding nie jest natywną cechą Neo4j -- cały graf
przechowywany jest na jednym węźle lub replikowany w całości na wszystkie
węzły klastra. Dla bardzo dużych grafów (miliardy węzłów i krawędzi) jest
to istotne ograniczenie. Skalowanie pionowe jest obsługiwane dobrze -- Neo4j
agresywnie cachuje często odwiedzane fragmenty grafu w pamięci RAM, co
przekłada się na bardzo niskie czasy odpowiedzi przy nawigacji po grafie.

\subsection{Bezpieczeństwo}

Neo4j oferuje kontrolę dostępu opartą na rolach (RBAC). Wbudowane role
obejmują uprawnienia do odczytu, zapisu i administracji, a użytkownik może
definiować własne role z granularnym dostępem do konkretnych baz danych,
etykiet węzłów i typów relacji. Uwierzytelnianie obsługuje hasła, certyfikaty
X.509 oraz integrację z LDAP i SSO przez OpenID Connect (Enterprise Edition).

Komunikacja między klientem a serwerem może być szyfrowana przez TLS.
Enterprise Edition oferuje też szyfrowanie danych w spoczynku oraz
zaawansowany audyt operacji. Od wersji 4.x Neo4j obsługuje wiele baz danych
w ramach jednej instancji (multi-database), co pozwala izolować dane różnych
aplikacji lub klientów na poziomie silnika.

\subsection{Awaryjność}

Trwałość danych w Neo4j zapewnia mechanizm Write-Ahead Log (WAL) -- każda
operacja jest najpierw zapisywana do dziennika transakcji, a dopiero potem
stosowana do plików grafu. W przypadku awarii serwer odtwarza stan
z dziennika przy kolejnym uruchomieniu. Neo4j obsługuje transakcje ACID
na poziomie pojedynczej operacji Cypher.

Causal Clustering (Enterprise Edition) zapewnia automatyczny failover --
jeśli węzeł primary przestanie odpowiadać, pozostałe węzły przeprowadzają
wybory lidera i jeden z secondary przejmuje rolę primary. Czas przełączenia
zależy od konfiguracji timeoutów, zazwyczaj wynosi kilka do kilkunastu
sekund. Kopie zapasowe można tworzyć przez \texttt{neo4j-admin backup}
(online, bez zatrzymywania serwera w Enterprise Edition) lub przez
zrzut bazy \texttt{neo4j-admin dump} (offline).

\subsection{Migracje}

Migracja danych do Neo4j realizowana jest najczęściej przez narzędzie
\texttt{neo4j-admin import} -- szybki importer wsadowy przyjmujący dane
w formacie CSV, zdolny do załadowania miliardów węzłów i krawędzi
w stosunkowo krótkim czasie. Alternatywą dla mniejszych zbiorów jest
użycie języka Cypher z instrukcją \texttt{LOAD CSV} lub bibliotek jak
APOC (Awesome Procedures on Cypher), które oferują dodatkowe procedury
importu z różnych formatów i źródeł.

Migracja z relacyjnych baz danych wymaga przemyślenia modelu danych --
tabele stają się etykietami węzłów, a klucze obce stają się relacjami.
Narzędzie \texttt{neo4j-etl-tool} (ETL Tool) wspiera automatyczną konwersję
schematu relacyjnego na model grafowy. Wersjonowanie schematu grafu nie ma
ustandaryzowanych narzędzi klasy Flyway czy Liquibase -- zarządzanie zmianami
schematu realizuje się zazwyczaj przez własne skrypty Cypher lub procedury APOC.

\subsection{Integracje}

Neo4j oferuje oficjalne sterowniki dla popularnych języków programowania:
Python (neo4j-driver), Java, JavaScript/Node.js, Go, .NET i Ruby, wszystkie
oparte na protokole Bolt. Dostępna jest też biblioteka OGM (Object Graph
Mapper) dla Javy i Spring Data Neo4j dla ekosystemu Spring.

APOC (Awesome Procedures on Cypher) to de facto standardowa biblioteka
rozszerzeń dla Neo4j, oferująca setki dodatkowych procedur i funkcji do
importu, eksportu, manipulacji grafem i integracji z zewnętrznymi systemami.
GDS (Graph Data Science Library) to oficjalna biblioteka algorytmów grafowych
-- PageRank, community detection, shortest path, link prediction i wiele
innych -- wykonywanych bezpośrednio w bazie. Neo4j AuraDB to zarządzana
wersja chmurowa dostępna na AWS, Google Cloud i Azure.

\subsection{Obszary biznesowe}

Neo4j najlepiej sprawdza się wszędzie tam, gdzie relacje między danymi są
równie ważne jak same dane i gdzie ich liczba oraz złożoność jest duża:

\begin{itemize}
    \item \textbf{Systemy rekomendacji} -- filtrowanie kolaboratywne
    i rekomendacje oparte na podobieństwie użytkowników i produktów są
    naturalnym zastosowaniem grafu. Stosowane przez Netflix, eBay i inne
    platformy streamingowe i e-commerce.
    \item \textbf{Wykrywanie oszustw} -- analiza powiązań między kontami,
    transakcjami i urządzeniami pozwala identyfikować podejrzane wzorce
    trudne do wykrycia w modelach relacyjnych. Używane przez banki
    i instytucje finansowe.
    \item \textbf{Grafy wiedzy i wyszukiwarki} -- przechowywanie encji
    i ich wzajemnych powiązań (knowledge graph) jest podstawą semantycznego
    wyszukiwania. Google Knowledge Graph jest przykładem tej klasy systemów.
    \item \textbf{Sieci społecznościowe} -- modelowanie użytkowników,
    znajomości, followersów i interakcji jako grafu jest naturalnym
    podejściem. LinkedIn używa grafowych baz danych do modelowania sieci
    zawodowej.
    \item \textbf{Zarządzanie siecią i infrastrukturą IT} -- modelowanie
    zależności między węzłami sieci, serwerami i usługami jako grafu
    ułatwia analizę wpływu awarii i planowanie zmian.
    \item \textbf{Systemy zgodności i regulacyjne} -- śledzenie łańcuchów
    własności, powiązań między podmiotami i przepływów finansowych
    w kontekście AML (Anti-Money Laundering) i KYC (Know Your Customer).
\end{itemize}

\subsection{Zalety}

\begin{itemize}
    \item \textbf{Natywny model grafowy} -- dane i relacje przechowywane
    są w strukturze optymalnej dla nawigacji po grafie. Przechodzenie
    przez relacje (graph traversal) ma złożoność zależną od lokalnego
    sąsiedztwa węzła, a nie od całkowitego rozmiaru bazy.
    \item \textbf{Język Cypher} -- deklaratywny, czytelny język zapytań
    zaprojektowany specjalnie dla grafów. Zapytania o wzorce w grafie
    (pattern matching) są znacznie bardziej zwięzłe niż ich odpowiedniki
    w SQL z wieloma JOIN-ami.
    \item \textbf{Wydajność przy głębokich powiązaniach} -- zapytania
    wymagające przejścia przez wiele poziomów relacji (np. znajomi znajomych,
    ścieżki w grafie) są w Neo4j wielokrotnie szybsze niż odpowiedniki
    w relacyjnych bazach danych.
    \item \textbf{Graph Data Science Library} -- bogata biblioteka
    gotowych algorytmów grafowych (PageRank, Louvain, betweenness centrality,
    node2vec) wykonywanych bezpośrednio w bazie bez eksportowania danych.
    \item \textbf{Intuicyjne modelowanie} -- model węzłów i relacji
    jest naturalnym odwzorowaniem rzeczywistości w wielu domenach,
    co ułatwia projektowanie i komunikację z ekspertami domenowymi.
    \item \textbf{APOC} -- rozbudowana biblioteka procedur rozszerzająca
    możliwości Cypher o setki dodatkowych funkcji, od manipulacji danymi
    po integrację z zewnętrznymi systemami.
\end{itemize}

\subsection{Wady i ograniczenia}

\begin{itemize}
    \item \textbf{Brak natywnego shardingu} -- cały graf musi mieścić się
    na jednym węźle lub być replikowany w całości, co ogranicza skalowalność
    przy bardzo dużych zbiorach danych.
    \item \textbf{Słabsza wydajność przy zapytaniach globalnych} -- operacje
    wymagające pełnego skanowania grafu (np. agregacje po wszystkich węzłach
    danego typu) są wolniejsze niż odpowiedniki w relacyjnych bazach danych
    z dobrze zaindeksowanymi tabelami.
    \item \textbf{Ograniczone narzędzia analityczne} -- Neo4j nie jest
    silnikiem OLAP i nie nadaje się do klasycznej analityki biznesowej
    opartej na raportach tabelarycznych. Integracja z narzędziami BI
    jest ograniczona w porównaniu do silników relacyjnych.
    \item \textbf{Mniejszy ekosystem} -- w porównaniu do PostgreSQL czy
    MySQL Neo4j ma mniej dostępnych sterowników, mniejszą społeczność
    i mniej gotowych rozwiązań, co może utrudniać znalezienie pomocy
    przy niestandardowych problemach.
    \item \textbf{Licencja Community Edition} -- wersja darmowa nie obsługuje
    klastrowania ani zaawansowanych funkcji bezpieczeństwa, co w praktyce
    wymusza zakup licencji Enterprise dla zastosowań produkcyjnych
    wymagających wysokiej dostępności.
    \item \textbf{Krzywa uczenia Cypher} -- mimo że Cypher jest czytelny,
    efektywne pisanie zapytań grafowych wymaga innego sposobu myślenia
    niż SQL, co może być barierą dla zespołów przyzwyczajonych do
    relacyjnych baz danych.
    \item \textbf{Narzut przy prostych operacjach CRUD} -- dla prostych
    operacji na pojedynczych rekordach bez złożonych powiązań Neo4j
    generuje większy narzut niż silniki relacyjne, co potwierdzają
    wyniki benchmarków przeprowadzonych w niniejszej pracy.
\end{itemize}

\subsection{Udogodnienia}

\begin{itemize}
    \item \textbf{Neo4j Browser} -- wbudowany interfejs webowy do
    interaktywnego odpytywania bazy w języku Cypher z wizualizacją
    wyników jako grafu. Pozwala eksplorować strukturę danych bez
    dodatkowych narzędzi.
    \item \textbf{Neo4j Bloom} -- narzędzie do wizualnej eksploracji
    grafu dla użytkowników biznesowych, bez konieczności znajomości
    Cypher. Pozwala wyszukiwać wzorce i przeglądać powiązania
    w intuicyjnym interfejsie graficznym.
    \item \textbf{EXPLAIN i PROFILE} -- Neo4j oferuje szczegółową analizę
    planów wykonania zapytań Cypher przez \texttt{EXPLAIN} (plan bez
    wykonania) i \texttt{PROFILE} (plan z rzeczywistymi statystykami),
    co pozwala diagnozować i optymalizować wolne zapytania.
    \item \textbf{APOC} -- biblioteka ponad 450 procedur i funkcji
    rozszerzających Cypher, obejmujących import/eksport danych, operacje
    na grafie, integracje z zewnętrznymi API i narzędzia do refaktoryzacji
    schematu.
    \item \textbf{Graph Data Science Playground} -- interaktywne środowisko
    do eksperymentowania z algorytmami grafowymi bezpośrednio w przeglądarce,
    bez konieczności pisania kodu.
    \item \textbf{Cypher jako standard} -- język Cypher jest standaryzowany
    przez openCypher i implementowany przez inne silniki grafowe (RedisGraph,
    Memgraph, AgensGraph), co ułatwia ewentualną migrację.
\end{itemize}